%% ==========================================================================
%%  FINAL REPORT --- SECTION 9, LIMITATIONS
%%  Each item is stated once, plainly, with where the evidence for it lives.
%% ==========================================================================
\section{Limitations}
\label{sec:limits}

We collect here everything a reader should discount the results by. Most of it
has been said at the point of use. It is repeated because a limitation stated
once, in the middle of a paragraph about something else, is easy to miss.

\paragraph{Semi-honest, with at most one corrupted server.} Every guarantee
assumes the servers follow the protocol. A deviating server can corrupt the
model and the answers, and nothing detects it. \textbf{Two colluding servers
learn everything}: any two hold all three shares of every value, and the two PIR
servers together reconstruct every fetched index. This is the assumption the
whole system rests on, and it is \nudgeplain{}'s own.

\paragraph{A sketch, not a proof.} Section~\ref{sec:threat:sketch} walks every
message and names a simulator, but there is no reduction and no hybrid argument.
The revealing normaliser is secure only against the weaker $\Fleaky$, and every
$\bits = 64$ quality number comes from it. What repeated training rounds reveal
about a user whose consumption feeds them is open.

\paragraph{Scale.} Every measurement is on MovieLens-100K. We derived the
truncation schedule at MovieLens-1M and Netflix scale, but did not run private
training there. A dense shared $6{,}040 \times 3{,}706$ matrix and its transpose,
for three parties in one process, need about 2.1\,GB, more than the development
machine could spare. The ring-width conclusions at larger scale are therefore
analytic.

\paragraph{Emulated networks, and one machine.} Wide-area results come from
channel-level emulation on a single laptop, not from servers in separate
locations or from kernel network emulation, neither of which the machine
supports. Training's wide-area times are computed from exact round and byte
counts, which is faithful for a round-bound protocol but ignores congestion and
jitter.

\paragraph{Absolute times are unreliable.} Thermal throttling moves absolute
timings by up to a factor of four between runs. We quote ratios, and where a
derived number moved between committed runs we give both values: the PIR
crossover link speed (8.0 and 22.9\,Gbit/s) and the denial-of-service
amplification (121 and 142 times). Anyone needing absolute numbers should rerun
on a machine that holds its clock.

\paragraph{One factor at a time.} The sweep varies each parameter around one
baseline. Interactions between parameters, for instance between $\dime$ and
$\bits$, are not measured.

\paragraph{Delivery is not batched.} Each recommendation costs a round trip, so
on a wide-area link the full catalogue download is 4.6 times faster than our
private read. Batching the $\topk$ queries into one trip is straightforward and
would reverse the result by our model, but it is not built.

\paragraph{The harvest runs as a library path, not in the deployed loop.} The
harvest, its conversion and the weight check are implemented and tested, but the
networked servers do not yet accumulate consumption across sessions or retrain
from it. Training in every experiment reads MovieLens ratings from disk. The
loop is closed in code and in tests, not in a long-running deployment.

\paragraph{Malicious clients are only partly handled.} The weight check stops a
client from inflating its vote. It does not stop a forged key from redistributing
weight, and it does nothing about the denial of service that full-domain
evaluation invites. Both need an audit in the style of \textsc{Sabre}, which we
did not build. The poisoning study places the collapse between eight and sixteen
colluders because only powers of two were run.

\paragraph{Leakage we do not analyse.} Timing on the wire, as opposed to server
compute time, is dominated by the operating system and was not studied. Query
counts and timing are leaked by design. The number of ratings each user
contributes is leaked unless every user submits the same number. And hiding the
fetch does nothing about what the recommendation reveals to anyone who later
observes the user act on it.

\paragraph{The server download of scores.} Because the user selects the top
$\topk$ locally, each server sends the user a share of all $\items$ scores. That
is cheap at 1{,}682 items and grows linearly with the catalogue, where an
oblivious selector run by the servers would send only $\topk$ indices at a large
cost in rounds.

\paragraph{No generic-MPC baseline.} The evaluation plan included running the
same factorisation in a general MPC framework. That needs Linux or Docker, which
the development machine lacks, so the comparison with generic MPC is absent.

\paragraph{Evaluation choices.} One dataset split, one seed for the attack, and
$\dime = 16$ throughout the leakage study. Our nDCG uses binary relevance at a
rating of four or more, so absolute values are comparable within this report and
not with \nudgeplain{}'s numbers.
